% !TeX root = ../../Book3.tex
% 纲要标题：## 附录 B　规范场、指标与特殊完整群
% 纲要位置：几何学计划纲要.md，第 3209 行。
% 正文状态：未撰写；本文件目前是框架。

% 纲要内容开始
% 
% > 利用正文丛联络和椭圆工具建立可计算的规范场模型；深层紧性和模空间理论精确引用。
% 
% * **Yang–Mills 变分**
%   * 联络空间及规范群
%   * 曲率能量和 Euler–Lagrange 方程
%   * 四维自对偶分解
%   * 拓扑下界与瞬子模型
% 
% * **椭圆化和模空间**
%   * Coulomb 规范的局部问题
%   * 形变复形与指标
%   * Uhlenbeck 紧性的准确输入
%   * 可去奇点及气泡能量
% 
% * **自旋与平行结构**
%   * Dirac 核和指标约束
%   * 平行自旋子与特殊完整群
%   * Kähler、超 Kähler 和例外模型
% 
% ---
% 
% 纲要内容结束

\chapter{规范场、指标与特殊完整群}
\label{g3:app:B}
\BookChapterNumber{g3:app:B}

待撰写

\section{Yang–Mills 变分}
\label{g3:sec:B01}
待撰写

\subsection{联络空间及规范群}
\label{g3:subsec:B01:01}
待撰写

\subsection{曲率能量和 Euler–Lagrange 方程}
\label{g3:subsec:B01:02}
待撰写

\subsection{四维自对偶分解}
\label{g3:subsec:B01:03}
待撰写

\subsection{拓扑下界与瞬子模型}
\label{g3:subsec:B01:04}
待撰写

\section{椭圆化和模空间}
\label{g3:sec:B02}
待撰写

\subsection{Coulomb 规范的局部问题}
\label{g3:subsec:B02:01}
待撰写

\subsection{形变复形与指标}
\label{g3:subsec:B02:02}
待撰写

\subsection{Uhlenbeck 紧性的准确输入}
\label{g3:subsec:B02:03}
待撰写

\subsection{可去奇点及气泡能量}
\label{g3:subsec:B02:04}
待撰写

\section{自旋与平行结构}
\label{g3:sec:B03}
待撰写

\subsection{Dirac 核和指标约束}
\label{g3:subsec:B03:01}
待撰写

\subsection{平行自旋子与特殊完整群}
\label{g3:subsec:B03:02}
待撰写

\subsection{Kähler、超 Kähler 和例外模型}
\label{g3:subsec:B03:03}
待撰写
